
\documentclass[11pt]{article}
\usepackage{amsmath}
\usepackage{geometry}
\usepackage{booktabs}
\usepackage{siunitx}
\usepackage{hyperref}
\geometry{margin=1in}

\title{Derivation of Pharmacokinetic Rate Constants from Volumes and Clearances}
\author{}
\date{}

\begin{document}
\maketitle

\section*{Introduction}
In compartmental pharmacokinetic models, drug distribution and elimination are described using rate constants that govern the movement of drug between compartments and out of the system. These rate constants are derived from physiological parameters such as compartment volumes and intercompartmental clearances.

\section*{Compartment Definitions}
We consider a three-compartment model:
\begin{itemize}
    \item Central compartment: volume $V_1$
    \item Peripheral compartment 1: volume $V_2$
    \item Peripheral compartment 2: volume $V_3$
\end{itemize}
Intercompartmental clearances:
\begin{itemize}
    \item $Q_2$: clearance between central and peripheral compartment 1
    \item $Q_3$: clearance between central and peripheral compartment 2
\end{itemize}
Systemic clearance from the central compartment is denoted as $Cl$.

\section*{Rate Constant Derivations}
The rate constants are defined as follows:

\subsection*{Elimination Rate Constant}
\[
k_{10} = \frac{Cl}{V_1}
\]
This represents the rate of drug elimination from the central compartment.

\subsection*{Intercompartmental Rate Constants}
\begin{align*}
k_{12} &= \frac{Q_2}{V_1} \quad \text{(central to peripheral 1)} \\
k_{21} &= \frac{Q_2}{V_2} \quad \text{(peripheral 1 to central)} \\
k_{13} &= \frac{Q_3}{V_1} \quad \text{(central to peripheral 2)} \\
k_{31} &= \frac{Q_3}{V_3} \quad \text{(peripheral 2 to central)}
\end{align*}

These equations assume that the intercompartmental clearance $Q$ is symmetric and that the rate constants are first-order with respect to drug concentration.

\section*{Units}
\begin{itemize}
    \item Clearance ($Cl$, $Q_2$, $Q_3$): \si{L/min} or \si{mL/kg/min}
    \item Volume ($V_1$, $V_2$, $V_3$): \si{L} or \si{L/kg}
    \item Rate constants ($k_{ij}$): \si{min^{-1}}
\end{itemize}

\section*{Conclusion}
By combining compartment volumes and intercompartmental clearances, we can derive the rate constants that define drug movement and elimination in a multi-compartment pharmacokinetic model. These constants are essential for simulating drug concentration-time profiles and for designing dosing regimens.

\end{document}
